\section{Conclusions and Outlook}\label{sec:conclusion}

In this work we have initiated a categorical study of UC for static systems, giving rigorous diagrammatic proofs of basic results such as the composition theorems. Our theory holds generally and gives a template for producing similar theories where ITMs are replaced by other computational systems. Moreover, categorical thinking fixed some prior oversights and resulted in several beneficial changes to UC (doing away with identities of machines altogether, allowing both the environment and the adversary to consist of multiple machines) while remaining inter-translatable with UC. 

Our work also opens up fertile grounds for future work:
\begin{itemize}
\item Techniques used in~\cite{BK23} can readily be used to develop other variants of the theory. For example, we can develop a version where honest machines are partitioned into $n$ sets ($n$ parties), with all communication across the partitions mediated by a shared functionality. We expect that the resulting categorical theory stands in a similar relationship to the simpler variant of UC in~\cite{DBLP:conf/crypto/CanettiCL15} as our theory here relates to UC for static systems. Similarly, we could readily reason about (abstract) distances up to $\varepsilon>0$, rather than up to indistinguishability. 

\item Can we model full UC similarly? Existing categorical tools are not well-suited for studying a dynamically changing network, making this a difficult question. 
  
\item As mentioned in the introduction, our theory is not an instance of~\cite{BK23} nor vice versa. This suggests  finding a unified categorical framework capturing both. However, we believe that in order to  avoid ``premature generalization'', one should first study categorically other approaches to composability such as
Abstract/Constructive Cryptography~\cite{MR11,Mau11} in general and random systems~\cite{M02} as a specific instance, quantum UC~\cite{Unr10}, reactive simulatability~\cite{PW01,BPW04,BPW07}, the IITM model~\cite{KTR20}, and programming-language-based approaches such as IPDL~\cite{DBLP:journals/pacmpl/GancherSFSM23}, ILC~\cite{LHM19}, state-separating proofs~\cite{BDF+18} and others~\cite{MT13,HS15}. 

\item The long-term goal is to investigate whether recasting various frameworks under a categorical umbrella would facilitate their comparison and pave the way for translating results across them. (Indeed, \cite[Theorem~4.10]{BK23} demonstrates an instance where symmetric monoidal functors preserve security.) This could be achieved via a web of security-preserving functors, thereby establishing when a protocol secure in one framework remains secure in others, and allowing for the derivation of a unified composition theorem.

\end{itemize}
\paragraph*{AI disclosure} This manuscript was proofread and edited with the assistance of large language models and subsequently revised by the authors, who are responsible for its content.






